% Abstract and keywords; heading structure follows the thesis template.
\begin{center}
{\Large\sffamily\bfseries Timing Synchronization for Orthogonal\\
Frequency-Division Multiplexing Systems\par}
\vspace{2mm}
{\large Zhuoyang Zhou\par}
\vspace{2mm}
{\large Statistics\par}
\vspace{2mm}
{\Large\bfseries Abstract\par}
\end{center}

This study compares four representative preamble-based timing estimators for orthogonal frequency-division multiplexing (OFDM): Schmidl--Cox (S\&C), Park-type, Minn-type and CAZAC-based (constant-amplitude zero-autocorrelation) methods. Two MATLAB Monte Carlo experiments each evaluate 880,000 estimates in additive white Gaussian noise (AWGN) and fixed multipath at signal-to-noise ratios of 0--20~dB. With equal training duration and energy, metrics include timing mean squared error (MSE), exact and cyclic prefix (CP) safety probabilities, and bit error rate (BER).

CAZAC achieves 219,997 exact estimates out of 220,000 with a known template. At 20~dB AWGN, S\&C achieves CP-safe and exact probabilities of 0.9978 and 0.0006; squared bias accounts for 97.36\% of its MSE. Park and Minn still exhibit occasional large errors. In the separate BER experiment at 20~dB AWGN, oracle phase correction using true timing error reduces S\&C bit errors from 639,428 to one out of 1,280,000 bits. Exact timing, CP-safe placement and reliable detection are thus distinct objectives.

\vspace{2mm}
\noindent\textbf{Keywords}: OFDM, timing synchronization, Monte Carlo simulation, mean squared error, cyclic prefix, bit error rate.
